
\documentclass[a4paper,11pt]{article}
\usepackage[a4paper, margin=8em]{geometry}

% usa i pacchetti per la scrittura in italiano
\usepackage[french,italian]{babel}
\usepackage[T1]{fontenc}
\usepackage[utf8]{inputenc}
\frenchspacing 

% usa i pacchetti per la formattazione matematica
\usepackage{amsmath, amssymb, amsthm, amsfonts}

% usa altri pacchetti
\usepackage{gensymb}
\usepackage{hyperref}
\usepackage{standalone}

% imposta il titolo
\title{Appunti Comunicazioni Numeriche}
\author{Luca Seggiani}
\date{2026}

% imposta lo stile
% usa helvetica
\usepackage[scaled]{helvet}
% usa palatino
\usepackage{palatino}
% usa un font monospazio guardabile
\usepackage{lmodern}

\renewcommand{\rmdefault}{ppl}
\renewcommand{\sfdefault}{phv}
\renewcommand{\ttdefault}{lmtt}

% disponi il titolo
\makeatletter
\renewcommand{\maketitle} {
	\begin{center} 
		\begin{minipage}[t]{.8\textwidth}
			\textsf{\huge\bfseries \@title} 
		\end{minipage}%
		\begin{minipage}[t]{.2\textwidth}
			\raggedleft \vspace{-1.65em}
			\textsf{\small \@author} \vfill
			\textsf{\small \@date}
		\end{minipage}
		\par
	\end{center}

	\thispagestyle{empty}
	\pagestyle{fancy}
}
\makeatother

% disponi teoremi
\usepackage{tcolorbox}
\newtcolorbox[auto counter, number within=section]{theorem}[2][]{%
	colback=blue!10, 
	colframe=blue!40!black, 
	sharp corners=northwest,
	fonttitle=\sffamily\bfseries, 
	title=~\thetcbcounter: #2, 
	#1
}

% disponi definizioni
\newtcolorbox[auto counter, number within=section]{definition}[2][]{%
	colback=red!10,
	colframe=red!40!black,
	sharp corners=northwest,
	fonttitle=\sffamily\bfseries,
	title=~\thetcbcounter: #2,
	#1
}

% U.D.M
\newcommand{\amp}{\ensuremath{\mathrm{A}}}
\newcommand{\volt}{\ensuremath{\mathrm{V}}}
\newcommand{\meter}{\ensuremath{\mathrm{m}}}
\newcommand{\second}{\ensuremath{\mathrm{s}}}
\newcommand{\farad}{\ensuremath{\mathrm{F}}}
\newcommand{\henry}{\ensuremath{\mathrm{H}}}
\newcommand{\siemens}{\ensuremath{\mathrm{S}}}

% circuiti
\usepackage{circuitikz}
\usetikzlibrary{babel}

% disegni
\usepackage{pgfplots}
\pgfplotsset{width=10cm,compat=1.9}

% disponi codice
\usepackage{listings}
\usepackage[table]{xcolor}

\lstdefinestyle{codestyle}{
		backgroundcolor=\color{black!5}, 
		commentstyle=\color{codegreen},
		keywordstyle=\bfseries\color{magenta},
		numberstyle=\sffamily\tiny\color{black!60},
		stringstyle=\color{green!50!black},
		basicstyle=\ttfamily\footnotesize,
		breakatwhitespace=false,         
		breaklines=true,                 
		captionpos=b,                    
		keepspaces=true,                 
		numbers=left,                    
		numbersep=5pt,                  
		showspaces=false,                
		showstringspaces=false,
		showtabs=false,                  
		tabsize=2
}

\lstdefinestyle{shellstyle}{
		backgroundcolor=\color{black!5}, 
		basicstyle=\ttfamily\footnotesize\color{black}, 
		commentstyle=\color{black}, 
		keywordstyle=\color{black},
		numberstyle=\color{black!5},
		stringstyle=\color{black}, 
		showspaces=false,
		showstringspaces=false, 
		showtabs=false, 
		tabsize=2, 
		numbers=none, 
		breaklines=true
}

\lstdefinelanguage{javascript}{
	keywords={typeof, new, true, false, catch, function, return, null, catch, switch, var, if, in, while, do, else, case, break},
	keywordstyle=\color{blue}\bfseries,
	ndkeywords={class, export, boolean, throw, implements, import, this},
	ndkeywordstyle=\color{darkgray}\bfseries,
	identifierstyle=\color{black},
	sensitive=false,
	comment=[l]{//},
	morecomment=[s]{/*}{*/},
	commentstyle=\color{purple}\ttfamily,
	stringstyle=\color{red}\ttfamily,
	morestring=[b]',
	morestring=[b]"
}

% disponi sezioni
\usepackage{titlesec}

\titleformat{\section}
	{\sffamily\Large\bfseries} 
	{\thesection}{1em}{} 
\titleformat{\subsection}
	{\sffamily\large\bfseries}   
	{\thesubsection}{1em}{} 
\titleformat{\subsubsection}
	{\sffamily\normalsize\bfseries} 
	{\thesubsubsection}{1em}{}

% disponi alberi
\usepackage{forest}

\forestset{
	rectstyle/.style={
		for tree={rectangle,draw,font=\large\sffamily}
	},
	roundstyle/.style={
		for tree={circle,draw,font=\large}
	}
}

% disponi algoritmi
\usepackage{algorithm}
\usepackage{algorithmic}
\makeatletter
\renewcommand{\ALG@name}{Algoritmo}
\makeatother

% disponi numeri di pagina
\usepackage{fancyhdr}
\fancyhf{} 
\fancyfoot[L]{\sffamily{\thepage}}

\makeatletter
\fancyhead[L]{\raisebox{1ex}[0pt][0pt]{\sffamily{\@title \ \@date}}} 
\fancyhead[R]{\raisebox{1ex}[0pt][0pt]{\sffamily{\@author}}}
\makeatother

\begin{document}
% sezione (data)
\section{Lezione del 12-03-26}

% stili pagina
\thispagestyle{empty}
\pagestyle{fancy}

% testo
\subsection{Probabilità condizionata}
Nella lezione 2 avevamo visto la definizione di \textit{probabilità} (definizione classica in 2.2, frequentistica in 2.3, e infine assiomatica attraverso gli assiomi di Kolmogorov in 2.4). Vediamo adesso la definizione di un concetto che ne deriva direttamente, cioè la \textbf{probabilità condizionata}:
\begin{definition}{Probabilità condizionata}
	Siano $a$ e $b$ due eventi, e $P$ una funzione di probabilità valida. Chiamiamo probabilità condizionata di $a$ dato $b$:
	$$
	P[a | b] = \frac{P(a \cap b)}{P(b)}
	$$
\end{definition}

Possiamo interpretare $P[a|b]$ come la probabilità dell'evento $a$, dato che l'evento $b$ si è verificato, cioè la probabilità condizionata cattura l'\textit{informazione parziale} che il verificarsi dell'evento $b$ fornisce sull'evento $a$.

Notiamo che quando si dà una nuova definizione di probabilità (come quella condizionata su un solito evento $b$) dobbiamo comunque verificare che gli assiomi di Kolmogorov risultino verificati. Quindi:
\begin{itemize}
	\item L'assioma di \textit{normalizzazione} vale se si considera $\Omega$ come evento $a$, cioè:
		$$
	P[\Omega | b] = \frac{P(\Omega \cap b)}{P(b)} = 1
		$$
		in quanto $\Omega \cap b$ non è altro che $b$ stesso;
	\item L'assioma di \textit{non negatività} è banale, dal fatto che $P[a|b]$ è rapporto di valori positivi;
	\item L'assioma di \textit{additività} si verifica per via algebrica, dati $a_1$, $a_2$ disgiunti:
		$$
		P[a_1 | b] + P[a_1 | b] = \frac{P(a_1 \cap b) + P(a_2 \cap b)}{P(b)} = \frac{P((a_1 \cup a_2) \cap b)}{P(b)} = P[a_1 \cup a_2 | b]
		$$
		notando che se $a_1$ e $a_2$ sono disgiunti, lo saranno per forza anche $a_1 \cap b$ e $a_2 \cap b$, permettendoci di applicare l'assioma di additività definito su $P$. 
\end{itemize}

\par\smallskip

\noindent
\textbf{\textsf{Interpretazione frequentista}}

\noindent
Vediamo velocemente cosa significa prendere la probabilità condizionata di $a$ dato $b$ se si prende come funzione $P$ quella frequentistica vista in 2.3:
$$
P(a) = \frac{n_A}{N}
$$
dati $N$ esperimenti.
$$
P[a | b] = \frac{P(a \cap b)}{P(b)} = \frac{ \frac{n_{a \cap b}}{N} }{ \frac{n_b}{N} } = \frac{n_{a \cap b}}{n_b}
$$


\par\smallskip

\noindent
\textbf{\textsf{Teoremi di probabilità condizionata}}

\noindent
Vediamo quindi alcuni teoremi riguardanti la probabilità condizionata:
\begin{enumerate}
\item \begin{theorem}{Probabilità condizionata di eventi disgiunti}
	Dati due eventi disgiunti $a$ e $b$, si ha:
	$$
	P[a | b] = 0
	$$
\end{theorem}
Questo è banale dalla definizione, visto che $a \cap b = \emptyset$ e quindi $P(a \cap b) = 0$. \hfill $\square$

\item \begin{theorem}{Probabilita condizionata con $a \subset b$}
	Dati due eventi che rispettano $a \subset b$, si ha:
	$$
	P[a | b] \geq P(a)
	$$
\end{theorem}
Questo si ha notando che se $a \subset b$, allora $a \cap b = a$ per cui:
$$
P[a | b] = \frac{P(a \cap b)}{P(b)} = \frac{P(a)}{P(b)} \geq P(b)
$$
visto che $P(b) \geq 0$ dall'assioma di non negatività. Questo, intuitivamente, significa che un evento che si verifica sempre assieme ad un secondo evento è più probabile dato il secondo evento. \hfill $\square$


\item \begin{theorem}{Probabilita condizionata con $b \subset a$}
	Dati due eventi che rispettano $b \subset a$, si ha:
	$$
	P[a | b] = 1 
	$$
\end{theorem}
Questo è il duale del precedente, e si dimostra vedendo che $b se \subset a$, allora $a \cap b = b$, per cui:
$$
P[a | b] = \frac{P(a \cap b)}{P(b)} = \frac{P(v)}{P(b)} = 1
$$
easyyy \hfill $\square$

\end{enumerate}

\subsubsection{Teorema della probabilità totale}
Vediamo un importante teorema basato sulla definizione di probabilità condizionata.
\begin{theorem}{Teorema della probabilità totale}
	Dato un evento $b$ e una partizione di $\Omega$ data dagli eventi $a_1, ..., a_K$, si ha:
$$
P(b) = \sum_{i = 1}^K P[b | a_i] P(a_i)
$$
\end{theorem}

Questo si dimostra dicendo che, se $a_1, ..., a_k$ sono una partizione (quindi disgiunti) si può prendere:
$$
b = b \cap \Omega = b \cap \bigcup_{i=1}^K a_i = \bigcup_{i=1}^K b \cap a_i
$$
Con tutti gli eventi $b \cap a_i$ sempre disgiunti. A questo punto basta applicare l'additività nel calcolo della probabilità:
$$
P(b) = P\left( \bigcup_{i=1}^K b \cap a_i \right) = \bigcup_{i=1}^K P(b \cap a_i) = P[b | a_i] P(a_i)
$$
che è la tesi. \hfill $\square$

In altri termini, il teorema della probabilità totale permette di calcolare la probabilità di un effetto (evento $b$) a partire dalle probabilità di tutte le sue possibili cause (eventi $a_i$) e dalle probabilità delle diverse combinazioni causa/effetto (probabilità condizionate $P[b|a_i]$).

\subsubsection{Teorema di Bayes}
Il teorema di Bayes ci permette di invertire la probabilità totale. In particolare:
\begin{theorem}{Teorema di Bayes}
Note $P[a|b]$, $P(a)$ e $P(b)$, vale che:
$$
P[b|a] = \frac{P[a | b]P(b)}{P(a)}
$$
\end{theorem}

Questo deriva dalla definizione di probabilità condizionata:
$$
P(a \cap b) = P[a | b] P(b) = P[b | a] P(a)
$$
per cui, eguagliando:
$$
P[b|a] = \frac{P[a | b]P(b)}{P(a)}
$$
che è la tesi. \hfill $\square$

\subsubsection{Eventi indipendenti}
Supponiamo che il verificarsi dell'evento $a$ non fornisca alcuna informazione sull'evento $b$, cioè:
$$
P[b | a] = P(b), \quad P[a | b] = P(a)
$$
dove la seconda equazione viene direttamente dal teorema di Bayes. 
Diciamo che due eventi $a$ e $b$ di questo tipo sono \textbf{eventi indipendenti}. Per questi vale:
\begin{theorem}{Eventi indipendenti}
	Dati $a$, $b$ eventi indipendenti, si ha:
	$$
	P(a \cap b) = P(a)P(b)
	$$
\end{theorem}
Questo è di facile dimostrazione dalle due equazioni sopra, in quanto:
$$
P[a | b] = P(a) = \frac{P(a \cap b)}{P(b)} \quad P(a \cap b) = P(a)P(b)
$$
che è la tesi. \hfill $\square$

Notiamo che l'opposto di questa affermazione è:
$$
P(a \cap b) = 0
$$
nel caso di eventi disgiunti, che era il primo teorema che avevamo dimostrato per la probabilità condizionata. 

\par\smallskip

\noindent
\textbf{\textsf{Interpretazione frequentista}}

\noindent
Vediamo cosa significa il teorema degli eventi indipendenti nell'interpretazione frequentistica della probabilità $P$. Se $a$ e $b$ sono indipendenti, vale:
$$
P[b | a] = P(b), \quad P[a | b] = P(a)
$$
e in generale conoscere $b$ non dà alcuna informazione riguardo al conoscere $a$, e viceversa.

Questo significherà che:
$$
\frac{n_{a \cap b}}{n_a} \approx \frac{n_b}{N}, \quad \frac{n_{a \cap b}}{n_b} \approx \frac{n_a}{N}
$$
cioè rispettivamente:
\begin{itemize}
	\item La frequenza di $b$, nella sequenza di $N$ ripetizioni dell’esperimento, approssima la frequenza di $b$ nella sequenza di $n_a$ ripetizioni in cui si è presentato $a$;
	\item La frequenza di $a$, nella sequenza di $N$ ripetizioni dell’esperimento, approssima la frequenza di $a$ nella sequenza di $n_b$ ripetizioni in cui si è presentato $b$.
\end{itemize}

Questo vale e diventa esatto nel limite ad infinito, cioè si può dire:
$$
\lim_{N \rightarrow \infty} \frac{n_b}{N} = \frac{n_{a \cap b}}{n_a}, \quad \lim_{N \rightarrow \infty} \frac{n_a}{N} = \frac{n_{a \cap b}}{n_b}
$$
sempre prendendo $n_a$ ed $n_b$ come sottoinsiemi di tutte le istanze di esperimnento $N$.

\subsection{Richiami di calcolo combinatorio}
Concludiamo facendo alcuni richiami di \textbf{calcolo combinatorio}, in quanto spesso questo torna utile nel campo della probabilità.

\begin{definition}{Principio fondamentale del calcolo combinatorio}
Se una procedura è composta di $N$ passi, il primo dei quali può essere svolto in $k_1$ modi diversi, il secondo in $k_2$ modi diversi, e via dicendo fino a $k_N$, allora il numero $M$ di modi in cui si può realizzare la procedura è:
$$
M = k_1 \cdot k_2 \cdot ... \cdot k_N = \prod_{i = 1}^N k_i
$$
\end{definition}

Nel caso per ogni passo si abbia lo stesso numero $k$ di modi, questo semplifica a:
$$
M = k^N
$$

\newpage

\noindent
\textbf{\textsf{Disposizioni semplici}}

\noindent
Con \textbf{disposizione} (semplice, senza ripetizioni) di $N$ oggetti presi $k$ a $k$ si definisce una sequenza ordinata di $k$ oggetti scelti tra gli $N$ assegnati. Questa si calcola come:
$$
D_{N. k} = N (N - 1) (N - 2) ... (N - k + 1) = \frac{N!}{(N - k)!}
$$

\par\smallskip

\noindent
\textbf{\textsf{Permutazioni semplici}}

\noindent
Con \textbf{permutazione} (semplice, senza ripetizioni) di $N$ oggetti si definisce una disposizione con $k = N$ degli oggetti. Questa si calcola come:
$$
P_N = D_{N, N} = N!
$$

\par\smallskip

\noindent
\textbf{\textsf{Combinazioni semplici}}

\noindent
Per \textbf{combinazione} (semplice, senza ripetizioni) di $N$ oggetti presi $k$ a $k$ si intende una sequenza non ordinata di $k$ oggetti scelti fra gli $N$ dati (ovvero, due sequenze formate dagli stessi elementi in ordine diverso sono indistinguibili). Questa si calcola come:
$$
C_{N, k} = \frac{D_{N, k}}{P_N} = \frac{N!}{(N - k)! k!} = \binom{N}{k}
$$

\par\smallskip

\noindent
\textbf{\textsf{Disposizioni con ripetizione}}

\noindent
Con \textbf{disposizione con ripetizione} di $N$ oggetti presi $k$ a $k$ si definisce una sequenza ordinata di $k$ oggetti scelti tra gli $N$ assegnati con possibilità di ripetere lo stesso oggetto più volte. Questa si calcola come:
$$
D^{(r)}_{N,k} = N^k
$$

\par\smallskip

\noindent
\textbf{\textsf{Permutazioni con ripetizione}}

\noindent
Con \textbf{permutazione con ripetizione} di $N$ oggetti, di cui $n_1$ uguali tra loro, $n_2$ uguali tra loro, \dots, $n_M$ uguali tra loro, con
$$
n_1 + n_2 + \dots + n_M = N
$$
si definisce una disposizione degli oggetti in cui gli elementi uguali sono indistinguibili. Il numero di permutazioni distinte è:
$$
P^{(r)}_N = \frac{P_N}{ \prod_{i = 1}^M P_{n_i} } = \frac{N!}{n_1! \, n_2! \dots n_M!}
$$

\par\smallskip

\noindent
\textbf{\textsf{Combinazioni con ripetizione}}

\noindent
Per \textbf{combinazione con ripetizione} di $N$ oggetti presi $k$ a $k$ si intende una scelta non ordinata di $k$ oggetti tra gli $N$ dati, con possibilità di ripetere lo stesso oggetto più volte. Questa si calcola come:
$$
C^{(r)}_{N,k} =
\binom{N + k - 1}{k}
=
\frac{(N + k - 1)!}{k!(N - 1)!}
$$


\subsubsection{Esperimento aleatorio composto}
Consideriamo ora come combinare i risultati di due esperimenti aleatori separati. Prendiamo gli spazi campione $\Omega_1$ e $\Omega_2$:
$$
\Omega_1 = \{ \xi_1, \xi_2, ... \xi_N \}, \quad
\Omega_2 = \{ \lambda_1, \lambda_2, ... \lambda_N \}
$$

Indichiamo con $a_1$ il generico evento del primo esperimento e con $a_2$ il generico evento del secondo esperimento. Inoltre prendiamo le funzioni di probabilità $P_1$ e $P_2$.

Possiamo quindi definire l'\textbf{esperimento composto} (o a volte \textit{esperimento congiunto}):
\begin{definition}{Esperimento composto}
	Un esperimento composto da due esperimenti di spazio campione $\Omega_1$, $\Omega_2$ (detti \textit{componenti}), è definito con:
\begin{itemize}
	\item Come \textit{risultati} le coppie ordinate $(\xi_i, \lambda_j)$ dei risultati degli esperimenti componenti;
	\item Come \textit{spazio campione} il prodotto cartesiano $\Omega_1 \times \Omega_2$ degli spazi campione;
	\item Come \textit{eventi} tutti i sottoinsiemi di $\Omega$ di tipo $a_1 \times a_2$.
\end{itemize}
\end{definition}

Per specificare completamente l’esperimento composto dobbiamo assegnare la sua funzione probabilità $P$. Come minimo, vorremo che questa rispetti:
$$
P(a_1 \times \Omega_2) = P_1(a_1), \quad
P(a_2 \times \Omega_1) = P_2(a_2)
$$
cioè la probabilità dell'evento $a_1$ dato un qualsiasi esito dell'esperimento 2 deve eguagliare la probabilità dell'evento $a_1$ nel solo esperimento 1 (e viceversa). 

\par\smallskip

\noindent
\textbf{\textsf{Esperimenti indipendenti}}

\noindent
Consideriamo un evento $a_1$ nello spazio campione $\Omega_1$ e un evento $a_2$ nello spazio campione $\Omega_2$: vogliamo calcolare la probabilità dell’evento:
$$
a = a_1 \times a_2
$$
nello spazio campione $\Omega$ dell'esperimento composto a partire dalla conoscenza di $P_1(a_1)$ e $P_2(a_2)$. Tale problema non ha, in generale, una soluzione: dalla conoscenza delle leggi di probabilità dei singoli esperimenti non è possibile ricavare la legge di probabilità dell’esperimento composto. 

Una importate eccezione è rappresentata dagli \textbf{esperimenti indipendenti}. Se scegliamo di imporre le proprietà della funzione di probabilità $P$ definita sull'esperimento composto, otteniamo che, stabilito:
$$
a_1 \times a_2 = (a_1 \times \Omega_2) \cap (a_2 \times \Omega_1)
$$
si ha che l'indipendenza è verificata quando:
$$
P(a) = P(a_1 \times a_2) = P((a_1 \times \Omega_2) \cap (a_2 \times \Omega_1)) = P(a_1 \times \Omega_2) P(\Omega_1 \times a_2) = P_1(a_1) P_2(a_2)
$$
ovvero quando concediamo il passaggio $P((a_1 \times \Omega_2) \cap (a_2 \times \Omega_1)) = P(a_1 \times \Omega_2) P(\Omega_1 \times a_2)$, cioè:
\begin{theorem}{Esperimenti indipendenti}
Imporre che 2 esperimenti, formato un esperimento composto, sono indipendenti, significa imporre che $a_1 \times \Omega_2$ e $\Omega_1 \times a_2$ rappresentano sempre eventi indipendenti.
\end{theorem}

\subsection{Prove ripetute}
Vediamo quindi la teoria delle \textbf{prove ripetute}, assumendo un esperimento con spazio campione $\Omega$ binario:
$$
\Omega = \{ a, \bar{a} \}
$$
da cui si ricava subito riguardo alla funzione di probabilità che:
$$
P(a) = p \implies P(\bar{a}) = 1 - p = q
$$

Quello che vogliamo fare è trovare la possibilità che l'evento $a$ si ripeta $k$ volte su $N$ prove, in qualsiasi ordine. In altre parole, vogliamo trovare la probabilità che si abbiano $a$ elementi positivi nel caso dell'esperimento composto $N$ volte, assumendo che l'esperimento in sé per sé sia \textit{indipendente}.

Prima di tutto, calcoliamo la probabilità che si verifichi un \textit{singolo} evento che rispetta le condizioni date:
$$
P(e) = p^k q^{N - k} = p^k (1 - p)^{N - k}
$$

A questo punto vorremo conoscere il numero $R$ di sequenze distinte che contengono $k$ eventi $a$, che equivale alle combinazioni di $N$ oggetti presi $k$ a $k$:
$$
C_{N, k} = \frac{D_{N, k}}{P_N} = \frac{N!}{(N - k)! k!} = \binom{N}{k}
$$

Ciò che ricaviamo è la cosiddetta \textbf{formula di Bernoulli}:
\begin{theorem}{Formula di Bernoulli}
La possibilità $P$ che l'evento $a$ si ripeta $k$ volte su $N$ prove sarà data da:
$$
P = \frac{N!}{(N - k)! k!} p^k (1 - p)^{N - k}
$$
\end{theorem}

\end{document}
